%Configuração do idioma usado no documento
\usepackage[brazilian]{babel}

%Definição padrao das margens das paginas
\usepackage[top=3cm,                    
bottom=2cm,
left=3cm,
right=2cm]{geometry}

%%%%% Configuração do espacamento padrão do documento - INICIO  %%%%%%%%%%%%%%%%%%%%%%%%
\usepackage{setspace}  % pacote usado para define o espacamento entre linhas do texto e de partes especificas
\onehalfspacing % definindo espaçamento padrao de 1,5 nas entrelinhas
%%%%% Configuração do espacamento padrão do documento - FIM  %%%%%%%%%%%%%%%%%%%%%%%%

% Usado para melhorias de alinhamento e justificação do texto
\usepackage{microtype}

% Usado Para permitir utilizar nomes de cores em diversos pacotes
\usepackage[table]{xcolor}

% Usado Para identar também intentar a primeira linha do primeiro parágrafo das subseções
\usepackage{indentfirst}

% Pacote da fonte escolhida: Opções de fontes em: https://ctan.org/topic/font
\usepackage{opensans}                    


%%%%% Removendo da contagem a capa e a ficha catalográfica - INICIO  %%%%%%%%%%%%%%%%%%%%%%%%
\addtocounter{page}{-2} 
%%%%% Removendo da contagem a capa e a ficha catalográfica - FIM  %%%%%%%%%%%%%%%%%%%%%%%%

% Usado nas referencias de equações, através do comando \eqref
\usepackage{amsmath}

% Necessario para incluir a ficha catalográfica, anexos e apêndices
\usepackage{pdfpages} 


% Usado para remover a quebra de linha nos recuos das folhas de rosto e de aprovação
\usepackage{hyphenat}                   

% Usado para justificar o texto nos recuos das folhas de rosto e de eprovação
\usepackage{ragged2e}                   

%%%%% Pacote e configurações da paginação - INICIO %%%%%%%%%%%%%%%%%%%%%%%%
\setlength{\headheight}{15.3pt} %Corrige warning do header do pacote fancyhdr
\usepackage{fancyhdr} % Pacote utilizado para alterar a posicao da paginação

%Redefine o style plain
\pagestyle{fancy}
\fancyhf{} % Limpa todos os cabeçalhos e rodapés padrão
\fancyhead[R]{\thepage} % Coloca o número da página no canto superior direito
\renewcommand{\headrulewidth}{0pt} % Remove a linha do cabeçalho (opcional)
\fancypagestyle{plain}{
	\fancyhf{} % Limpa cabeçalhos e rodapés
	\fancyhead[R]{\thepage} % Número da página no canto superior direito
	\renewcommand{\headrulewidth}{0pt} % Remove a linha do cabeçalho
}
%%%%% Pacote e configurações da paginação - INICIO %%%%%%%%%%%%%%%%%%%%%%%%



%%%%% Configurações de links das referencias internas e externos - INICIO  %%%%%%%%%%%%%%%%%%%%%%%%
\usepackage{hyperref} % Cria os links das referencias internas e para links externos tb
\makeatletter
\hypersetup{
	pdftitle={\@title},
	pdfpagemode=FullScreen,
	colorlinks=true,      % Ativa a coloração dos links
	linkcolor=black,     % Cor dos links internos
	citecolor=black,     % Cor das citações
	urlcolor=black       % Cor dos links de URL
}
\makeatother
%%%%% Configurações de links das referencias internas e externos - FIM  %%%%%%%%%%%%%%%%%%%%%%%%



%%%%% Configuração do estilo CSL usado na bibilografia - INICIO %%%%%%%%%%%%%%%%%%%%%%%%%%%%%%%%
\usepackage[style=config/template-ibict-completo-20251112]{citation-style-language}
\addbibresource{3-pos-textuais/referencias.bib}
%%%%% Configuração do estilo CSL usado na bibilografia - FIM %%%%%%%%%%%%%%%%%%%%%%%%%%%%%%%%%%%


%%%%% Configurações do texto de Fonte das imagens - INICIO  %%%%%%%%%%%%%%%%%%%%%%%%
\usepackage{copyrightbox} % Necessario para adicionar origem das imagem e das tabelas
\makeatletter
\renewcommand{\CRB@setcopyrightfont}{%
	\color{black}
}
\makeatother
%%%%% Configurações do texto de Fonte das imagens - FIM  %%%%%%%%%%%%%%%%%%%%%%%%



%%%%% Evita reiniciar numeracoes das equações nos capitulos - INICIO  %%%%%%%%%%%%%%%%%%%%%%%%
\usepackage{chngcntr} % Permite modificar a numeração dos contadores: Equacoes, Tabelas, e Figuras
\counterwithout{equation}{chapter} % Remove a reinicialização da numeração a cada capítulo
% A configuração de tabela e figura estão em  seus respectivos arquivos de listagem 
% juntamente com as suas demais configurações 
%%%%% Evita reiniciar numeracoes das equações nos capitulos - FIM  %%%%%%%%%%%%%%%%%%%%%%%%


%%%%% Pacote que permite criar indices - INICIO  %%%%%%%%%%%%%%%%%%%%%%%%
\usepackage{imakeidx} % Pacote para criar o índice
\makeindex[columns=3, options= -s config/estilo-indice.ist] % Configura quantidade de colunas do índice e o arquivo de estilo
%%%%% Pacote que permite criar indices - FIM  %%%%%%%%%%%%%%%%%%%%%%%%


% Permite alterar o formato das listas enumeradas 
\usepackage{enumitem}                   

%\usepackage{url}                        % Usado para formatar URLs no texto





%%%%% Formatação da lista de figura e de tabelas - INICIO  %%%%%%%%%%%%%%%%%%%%%%%%
\usepackage{titletoc} % Pacote para formatar sumário, lista de figuras e tabelas

%a configuração só funciona se for antes do documentclass

%Lista de Tabelas
\titlecontents{figure}[0pt]
{}
{\figurename\ \thecontentslabel\ - }
{}
{\titlerule*[0.4pc]{.}\contentspage}
[\addvspace{5pt}]

%Lista de Tabelas
\titlecontents{table}[0pt]
{}
{\tablename\ \thecontentslabel\ - }
{ }
{\titlerule*[0.4pc]{.}\contentspage}
[\addvspace{5pt}]
%%%%% Formatação da lista de figura e de tabelas - FIM  %%%%%%%%%%%%%%%%%%%%%%%%



%%%%% Formatação dos titulos dos capitulos, secao, subsecao, etc - INICIO  %%%%%%%%%%%%%%%%%%%%%%%%
\usepackage{titlesec}

\titlespacing{\chapter}
  {0pt}   % Espaço à esquerda
  {-17pt}   % Espaço acima do título do capítulo (é aqui que se reduz)
  {15pt}  % Espaço abaixo do título do capítulo

\titleformat{\chapter}[hang] % shape
{\normalsize\bfseries\uppercase} % format
{\thechapter}
{10pt}
{}

\titleformat{\section}[hang] % shape
{\normalsize\normalfont\uppercase} % format
{\thesection}
{10pt}
{}

\titleformat{\subsection}[hang] % shape
{\normalsize\bfseries} % format
{\thesubsection}
{10pt}
{}

\titleformat{\subsubsection}[hang] % shape
{\normalfont\normalsize} % format
{\thesubsubsection}
{10pt}
{}

\titleformat{\paragraph}[hang] % shape
{\normalfont\normalsize} % format
{\theparagraph}
{10pt}
{}

\titleformat{\subparagraph}[hang] % shape
{\normalfont\normalsize} % format
{\thesubparagraph}
{}
{}
%%%%% Formatação dos titulos dos capitulos, secao, subsecao, etc - FIM  %%%%%%%%%%%%%%%%%%%%%%%%

% pacote que gera textos aletorios, pode remover quando for realizar o trabalho
\usepackage{lipsum} 




%%%%% Definindo melhor configuração das imagens - INICIO %%%%%%%%%%%%%%%%%%%%%%%%
\usepackage{float}
%%%%% Definindo melhor configuração das imagens - FIM %%%%%%%%%%%%%%%%%%%%%%%%



%%%%% Definindo diretorios padrao do template e do projeto - INICIO %%%%%%%%%%%%%%%%%%%%%%%%
\usepackage{graphicx}		            % Inclusão de gráficos, pdfs
\graphicspath{{figuras/}{figuras/modelo}}
%%%%% Definindo diretorios padrao do template e do projeto - INICIO %%%%%%%%%%%%%%%%%%%%%%%%